\section*{第 \ref{ch:b3:lagrangian-mechanics} 章——拉格朗日力学}

\begin{solution}{exo:b3:lagrangian-mechanics:1}
(a) 2（碗面上的两个角度）。(b) 1：沿
斜面方向的坐标 $x$，转角由滚动与之相联系，$x =
R\phi$。(c) 3：$X$、$\theta_1$、$\theta_2$。(d) 5：
质心三个，杆的方向两个。(e) 2：每个各一个角度。
滚动条件是\emph{速度}之间的关系，$\dot x =
R\dot\phi$；在此平面问题中可直接积分得 $x = R\phi +
\text{常数}$，因而是完整的。（对于在平面上滚动的球，它
不可积，拉格朗日力学需要推广。）
\end{solution}

\begin{solution}{exo:b3:lagrangian-mechanics:2}
(a) $[L] = \unit{J}$；$[p_\theta] = \unit{kg.m^2/s}$——
是角动量，因为 $\theta$ 无量纲。(b) $p_\theta =
m\ell^2\dot\theta$ 为摆球绕支点的角动量。
(c) 由 $\ddot\theta = -(g/\ell)\theta$ 得 $T = 2\pi\sqrt{\ell/g}$。(d)
以平衡位置为零点度量 $E_p$，简谐振动的时间平均动能与势能相等，
故 $\langle
L\rangle = 0$，且整个周期内 $S = 0$——计算
证实了这一点：$S = \int_0^T\tfrac12 m\ell^2\theta_0^2[\omega^2\sin^2
\omega t - (g/\ell)\cos^2\omega t]\,\dd t = 0$，因为 $\omega^2 =
g/\ell$。
\end{solution}

\begin{solution}{exo:b3:lagrangian-mechanics:3}
(a) 设 $x$ 为 $m_1$ 的下降距离（于是 $m_2$ 上升 $x$）：$L =
\tfrac12(m_1 + m_2)\dot x^2 + (m_1 - m_2)gx$。(b) $(m_1 + m_2)\ddot x =
(m_1 - m_2)g$，故 $a = (m_1 - m_2)g/(m_1 + m_2)$。(c) 张力作用于
不可伸长细绳的两端，在任何允许的位移中其
功相互抵消，因此它从不进入 $L$。(d) 仅对 $m_1$ 应用牛顿定律：$T =
m_1(g - a) = 2m_1m_2g/(m_1 + m_2)$。
\end{solution}

\begin{solution}{exo:b3:lagrangian-mechanics:4}
(a) $\varphi$ 与 $z$：$p_\varphi = mr^2\dot\varphi$（
绕轴的角动量）与 $p_z = m\dot z$ 守恒。(b) $r$
在 $L$ 中通过 $r^2\dot\varphi^2$ 出现，因此它不是循环坐标：$\dot p_r =
mr\dot\varphi^2 \neq 0$。这并无问题——$p_r = m\dot r$ 为
常矢量 $\vect p$ 的径向分量，沿
\emph{转动}方向的分量不必为常数。(c) $m\ddot r =
mr\dot\varphi^2$：离心项，即使用
固连于转动方向的坐标所付出的代价。(d) 对距离为 $b$
的直线、以速率 $v$ 掠过时：$r = \sqrt{b^2 + v^2t^2}$，$r^2\dot\varphi = bv$；
于是 $\ddot r = b^2v^2/r^3 = r(bv/r^2)^2 = r\dot\varphi^2$。
\end{solution}

\begin{solution}{exo:b3:lagrangian-mechanics:5}
(a) 滑块位置 $(X + s\cos\alpha,\ -s\sin\alpha)$，故
\[
L = \tfrac12 M\dot X^2 + \tfrac12 m\big(\dot X^2 +
2\dot X\dot s\cos\alpha + \dot s^2\big) + mgs\sin\alpha .
\]
(b) $X$ 为循环坐标：$P = (M + m)\dot X + m\dot s\cos\alpha$
守恒——总水平动量，因为没有外部水平力
作用。(c) $s$ 方程为 $\ddot s + \ddot
X\cos\alpha = g\sin\alpha$；由 (b) 得 $\ddot X = -m\ddot s\cos\alpha/(M +
m)$，
\[
\ddot s = \frac{(M + m)g\sin\alpha}{M + m\sin^2\alpha} , \qquad
\ddot X = -\frac{mg\sin\alpha\cos\alpha}{M + m\sin^2\alpha} .
\]
(d) $M \to \infty$：$\ddot s \to g\sin\alpha$，即固定斜面；
$\alpha \to 90^\circ$：$\ddot s \to g$，$\ddot X \to 0$——自由落体
沿竖直面下落，楔块未受推动。
\end{solution}

\begin{solution}{exo:b3:lagrangian-mechanics:6}
(a) $L = \tfrac12 m\ell^2(\dot\theta^2 + \sin^2\theta\,\dot\varphi^2) +
mg\ell\cos\theta$。(b) $\varphi$ 为循环坐标：$p_\varphi =
m\ell^2\sin^2\theta\,\dot\varphi$，即竖直方向的角动量。(c)
消去 $\dot\varphi$ 后，$\tfrac12 m\ell^2\dot\theta^2 +
U_{\text{有效}}(\theta)$ 守恒，其中
\[
U_{\text{有效}}(\theta) = \frac{p_\varphi^2}{2m\ell^2\sin^2\theta}
 - mg\ell\cos\theta :
\]
在 $\theta = 0$ 与 $\theta = \pi$ 处为壁（$p_\varphi \neq 0$ 时），
二者之间有一个极小值——摆球在两个圆之间做章动。(d) 在
极小值处，$\dot\theta = 0$ 且 $\theta = \theta_0$ 为常数：
$p_\varphi^2\cos\theta_0/m\ell^2\sin^3\theta_0 = mg\ell\sin\theta_0$；
代入 $p_\varphi = m\ell^2\sin^2\theta_0\,\omega$ 得
$\ell\omega^2\cos\theta_0 = g$。
\end{solution}

\begin{solution}{exo:b3:lagrangian-mechanics:7}
(a) 摆球位于 $(X + \ell\sin\theta,\ -\ell\cos\theta)$：
\[
L = \tfrac12(M + m)\dot X^2 + m\ell\cos\theta\,\dot X\dot\theta
 + \tfrac12 m\ell^2\dot\theta^2 + mg\ell\cos\theta .
\]
(b) $X$ 为循环坐标：$(M + m)\dot X + m\ell\cos\theta\,\dot\theta$
守恒——整个系统的水平动量（导轨
只施加竖直方向的力）。(c) 小角度近似：$(M + m)\ddot X +
m\ell\ddot\theta = 0$ 与 $\ell\ddot\theta + \ddot X + g\theta = 0$；
消去 $\ddot X$，得 $\ell\ddot\theta\,[1 - m/(M + m)] = -g\theta$，
故 $\Omega^2 = (M + m)g/M\ell = (1 + m/M)\,g/\ell$。(d) $M \to \infty$：
支点固定，$\Omega^2 = g/\ell$。$M$ 很小时：小车
向与摆球相反的方向反冲，支点至摆球的摆动以二者之间的一点为中心
（固定的质心）发生，使有效摆长变短
——因此频率更高，且当 $M \to 0$ 时趋于发散。
\end{solution}

\begin{solution}{exo:b3:lagrangian-mechanics:8}
(a) $\vect A = \tfrac12 B(-y, x, 0)$：
$L = \tfrac12 m(\dot x^2 + \dot y^2 + \dot z^2) + \tfrac12 qB(x\dot y -
y\dot x)$。(b) $x$ 方程：$\dd(m\dot x - \tfrac12 qBy)/\dd t =
\tfrac12 qB\dot y$，即 $m\ddot x = qB\dot y$；同理 $m\ddot y =
-qB\dot x$：以 $\omega_{\text{c}} = qB/m$ 做圆周运动，$\dot z$
为常数。(c) $p_x = m\dot x - \tfrac12 qBy \neq m\dot x$；$x$ 与
$y$ 均不是循环坐标（$\partial L/\partial x = \tfrac12 qB\dot y$），因此
两个动量都不守恒——只有 $m\dot x -
qBy$ 这类组合守恒（验证：其导数为零）。(d) 在
柱坐标中 $A_\varphi = \tfrac12 Br$：
$L = \tfrac12 m(\dot r^2 + r^2\dot\varphi^2 + \dot z^2) + \tfrac12
qBr^2\dot\varphi$；$\varphi$ 为循环坐标，$p_\varphi = mr^2\dot\varphi +
\tfrac12 qBr^2$。在圆心位于轴上、半径为 $R$ 的圆上，
$\dot\varphi = -\omega_{\text{c}}$，故 $p_\varphi = -qBR^2 + \tfrac12
qBR^2 = -\tfrac12 qBR^2$。
\end{solution}

\begin{solution}{exo:b3:lagrangian-mechanics:9}
(a) $h = \dot\theta\,\partial L/\partial\dot\theta - L = \tfrac12
mR^2\dot\theta^2 - \tfrac12 m\omega^2R^2\sin^2\theta - mgR\cos\theta =
\tfrac12 mR^2\dot\theta^2 + U_{\text{eff}}$；$\dd h/\dd t =
\dot\theta\,[mR^2\ddot\theta + U_{\text{eff}}'(\theta)] = 0$，由
运动方程。(b) $E = \tfrac12 mR^2\dot\theta^2 + \tfrac12
m\omega^2R^2\sin^2\theta - mgR\cos\theta = h +
m\omega^2R^2\sin^2\theta$。(c) $P = \dd E/\dd t =
m\omega^2R^2\sin2\theta\,\dot\theta$：珠子向外
远离转轴时为正（电动机克服珠子的惯性做功），
返回时为负。(d) $U_{\text{eff}}'' = mgR\cos\theta -
m\omega^2R^2\cos2\theta$；令 $\cos\theta_{\text{平衡}} = g/\omega^2R$
此即 $m\omega^2R^2\sin^2\theta_{\text{平衡}}$，故
$\omega_{\text{振荡}} = \omega\sin\theta_{\text{平衡}} = \omega\sqrt{1 -
(g/\omega^2R)^2} \to 0$，当 $\omega^2 \to g/R$ 时，回复力
恰在两势阱合并时变得平缓。
\end{solution}

\begin{solution}{exo:b3:lagrangian-mechanics:10}
(a) 位置 $x_2 = \ell(\sin\theta_1 + \sin\theta_2)$ 等；保留
二次项，交叉速度项为
$m\ell^2\dot\theta_1\dot\theta_2$，即得所给 $L$。(b)
$2\ddot\theta_1 + \ddot\theta_2 = -2\omega_0^2\theta_1$ 和
$\ddot\theta_1 + \ddot\theta_2 = -\omega_0^2\theta_2$，其中 $\omega_0^2
= g/\ell$。(c) 代入 $\theta_i = a_i\cos\Omega t$：
$\det\begin{pmatrix} 2\omega_0^2 - 2\Omega^2 & -\Omega^2\\ -\Omega^2 &
\omega_0^2 - \Omega^2\end{pmatrix} = \Omega^4 - 4\omega_0^2\Omega^2 +
2\omega_0^4 = 0$，故 $\Omega_\pm^2 = (2 \mp \sqrt2)\omega_0^2$；
由第二行可得 $a_2/a_1 = \Omega^2/(\omega_0^2 - \Omega^2) =
\pm\sqrt2$：两摆球同向（慢模式），两摆球反向（快模式）。(d)
大振幅时，$\sin$ 与 $\cos$ 的耦合使方程
成为非线性；解不再满足叠加原理，相邻的初始
条件指数分离（混沌）。能量仍
严格守恒——$L$ 不显含时间——混沌关乎
可预测性，而非守恒。
\end{solution}

\begin{solution}{exo:b3:lagrangian-mechanics:11}
(a) $v = \sqrt{2gy}$ 与 $\dd s = \sqrt{1 + y'^2}\,\dd x$ 给出
泛函。(b) 仿照
\cref{prop:b3:lagrangian-mechanics:energy} 的计算，以 $x$ 代替时间：$\dd
h/\dd x = -\partial F/\partial x = 0$。(c) $h = -1/\sqrt{2gy(1 +
y'^2)}$，故 $y(1 + y'^2) = 2r$。对摆线，$y' =
\sin\phi/(1 - \cos\phi)$ 且 $1 + y'^2 = 2/(1 - \cos\phi)$，因此
$y(1 + y'^2) = 2r$。(d) $\dd t = \dd s/v = \sqrt{r/g}\,\dd\phi$（所有
$\phi$ 的依赖项相消），因此到达最低点（$\phi = \pi$）所需
$\pi\sqrt{r/g}$——与起点无关：摆线也是
等时曲线。通往 $(\pi r, 2r)$ 的直线滑道所需时间为 $\sqrt{\pi^2
+ 4}\,\sqrt{r/g} \approx 3.72\sqrt{r/g}$，约长 $18\%$（以
$\pi\sqrt{r/g}$ 为基准）。
\end{solution}

\begin{solution}{exo:b3:lagrangian-mechanics:12}
(a) $\dd t = \dd s/(c/n)$。(b) $h = -n(y)/\sqrt{1 + y'^2}$ 守恒；
$1/\sqrt{1 + y'^2} = \cos\theta$（$\theta$ 为斜率角）
$= \sin i$，其中 $i$ 自竖直方向量起：$n\sin i = \text{常数}$——
斯涅耳定律的连续应用。(c) 对近水平的光线
$n\cos\theta \approx \text{常数}$，当 $\theta$ 很小时给出
$\theta\,\dd\theta = \dd n/n \approx \beta\,\dd y$；由于 $\dd y =
\theta\,\dd x$，曲率为 $\dd\theta/\dd x = \beta$，即 $R =
1/\beta \approx \qty{83}{km}$，光线向\emph{上}弯曲（朝向更大的
$n$）。(d) 自人眼发出的光线，经半径为
$R$ 的弯曲后与路面相切，切点距离 $d = \sqrt{2hR} = \sqrt{2 \times 1.2 \times
8.3e4} \approx \qty{450}{m}$：超过该距离，路面本身
不可见——人看到的是向上折射的天空，即闪烁的“水”。
\end{solution}

\begin{solution}{pb:b3:lagrangian-mechanics:1}
\textbf{1.} $L = \tfrac12 m\ell^2\dot\theta^2 + mg\ell\cos\theta$；
$\ddot\theta = -(g/\ell)\sin\theta$。
\textbf{2.} $f_0 = \sqrt{g/\ell}/2\pi = \qty{0.79}{Hz}$。
\textbf{3.} 悬挂点固定，故 $h = E = \tfrac12
m\ell^2\dot\theta^2 - mg\ell\cos\theta$；由最低点到最高点，$\tfrac12
mv^2 = 2mg\ell$：$v = 2\sqrt{g\ell} = \qty{4.0}{m/s}$。
\textbf{4.} $\ddot\epsilon = +(g/\ell)\epsilon$：$\epsilon \propto
\eu^{t/\tau}$，其中 $\tau = \sqrt{\ell/g} = \qty{0.20}{s}$。
\textbf{5.} 倒立平衡位置存在，但
是指数不稳定的：任何偏斜，无论多小，每
\qty{0.2}{s} 便放大为原来的 $e$ 倍。
\textbf{6.} 摆球位于 $(\ell\sin\theta,\ a\cos\Omega t - \ell\cos\theta)$；
求导并平方：得所述 $\vect v^{\,2}$，其中交叉
项为 $-2a\Omega\ell\sin(\Omega t)\sin\theta\,\dot\theta$。
\textbf{7.} 若 $L' = L + \dd F(q,t)/\dd t$，则作用量改变
$F(q_2, t_2) - F(q_1, t_1)$，在固定端点的变分下为常数，
故驻定路径相同，方程也相同。
\textbf{8.} 交叉项为 $-ma\Omega\ell\sin(\Omega t)\sin\theta\,
\dot\theta = \dd[ma\Omega\ell\sin(\Omega t)\cos\theta]/\dd t -
ma\Omega^2\ell\cos(\Omega t)\cos\theta$；略去全导数项
以及仅含 $t$ 的项（$\tfrac12 ma^2\Omega^2\sin^2\Omega t$ 与
$-mga\cos\Omega t$），即得所述 $L$。
\textbf{9.} $m\ell^2\ddot\theta = -mg\ell\sin\theta +
ma\Omega^2\ell\cos(\Omega t)\sin\theta$。由于 $\ddot y_{\text{s}} =
-a\Omega^2\cos\Omega t$，故 $\ddot\theta =
-[(g + \ddot y_{\text{s}})/\ell]\sin\theta$：在支点的参考系中，
表观重力振荡。
\textbf{10.} $L$ 现在显含 $t$：$h$ 不守恒，
$E$ 亦不守恒——振动台通过支点向系统输入与输出能量。
\textbf{11.} $a/\ell = 0.05$；$\Omega = \qty{251}{rad/s}$，与
$\omega_0 = \qty{4.9}{rad/s}$ 相比，比值为 $51$。在一个驱动周期
（\qty{25}{ms}）内，重力几乎不改变 $\dot\theta$：摆球太
迟缓，无法跟随，只能微微颤动。
\textbf{12.} $\ddot\xi$ 是最大的导数（$\propto\Omega^2$），
驱动项是最大的力项：$\ddot\xi =
(a\Omega^2/\ell)\cos(\Omega t)\sin\Theta$。
\textbf{13.} 在固定 $\Theta$ 下积分两次：$\xi =
-(a/\ell)\cos(\Omega t)\sin\Theta$，振幅至多为 $a/\ell =
0.05$，即两度的微颤。
\textbf{14.} $\sin\theta \approx \sin\Theta + \xi\cos\Theta$，于是
\[
\ddot\Theta + \ddot\xi = -\frac{g}{\ell}(\sin\Theta + \xi\cos\Theta)
 + \frac{a\Omega^2}{\ell}\cos\Omega t\,(\sin\Theta + \xi\cos\Theta) .
\]
\textbf{15.} 取平均后，$\ddot\xi$、$\langle\cos\Omega t\rangle$
与 $\langle\xi\rangle$ 均消去；保留的交叉项为
$(a\Omega^2/\ell)\cos\Theta\,\langle\xi\cos\Omega t\rangle =
-(a^2\Omega^2/2\ell^2)\sin\Theta\cos\Theta$，由此得到所述的
$\Theta$ 的方程。
\textbf{16.} $\ddot\Theta = -(1/m\ell^2)\,U_{\text{eff}}'(\Theta)$，其中
$U_{\text{eff}} = mg\ell[-\cos\Theta + (a^2\Omega^2/4g\ell)
\sin^2\Theta]$——求导即可检验。
\textbf{17.} 与圆环的 $\sin^2$ 项相同，但
\emph{相反}符号：旋转在两侧挖出势阱，只能
使底部变平；竖直振动使底部势阱变陡，并在
\emph{顶部}处挖出一个新的势阱。
\textbf{18.} 慢驱动（$a^2\Omega^2 < 2g\ell$）：极小值位于 $\Theta =
0$，极大值位于 $\pi$——并无新情况。快驱动：极小值位于 $0$
\emph{与} $\pi$，两者之间由位于 $\cos\Theta^* =
-2g\ell/a^2\Omega^2$ 的极大值分隔；悬挂摆与直立摆均会
稳定振荡。
\textbf{19.} 在 $\pi$ 附近，取 $\Theta = \pi + \epsilon$：
$\ddot\epsilon = [g/\ell - a^2\Omega^2/2\ell^2]\,\epsilon$；稳定性
要求括号内的量为负，即 $a^2\Omega^2 > 2g\ell$。
\textbf{20.} $\Omega_{\text{c}} = \sqrt{2g\ell}/a = \qty{140}{rad/s}$：
$f_{\text{c}} = \qty{22}{Hz}$；峰值速率 $a\Omega_{\text{c}} =
\qty{2.8}{m/s}$，峰值加速度 $a\Omega_{\text{c}}^2 =
\qty{392}{m/s^2} \approx 40g$。
\textbf{21.} $\omega_{\text{慢}} = \sqrt{a^2\Omega^2/2\ell^2 -
g/\ell}$；当 $\Omega = 2\Omega_{\text{c}}$ 时，$a^2\Omega^2 = 8g\ell$，故
$\omega_{\text{慢}} = \sqrt{3g/\ell} = \qty{8.6}{rad/s}$：
\qty{1.4}{Hz}，比 \qty{45}{Hz} 的驱动慢三十倍。
\textbf{22.} $\xi_{\max} = (a/\ell)\sin\Theta$：若摆杆相对竖直方向倾斜 $10^\circ$，
则 $\xi_{\max} = 0.05\sin170^\circ = \qty{8.7}{mrad}
\approx 0.5^\circ$——普通照片中，扫帚看起来
静止不动。
\textbf{23.} 直立势阱达到两侧的极大值：
$|\Theta - \pi| < \arccos(2g\ell/a^2\Omega^2)$；当 $\Omega =
2\Omega_{\text{c}}$ 时，$\arccos\tfrac14 = 75^\circ$——一个相当
宽容的势阱。
\textbf{24.} 杂耍者利用反馈：眼睛测量倾斜，手
向侧面加速以抵消它。卡皮察的稳定化是
开环的——驱动从不知道摆在何处。
\textbf{25.} 以 \qty{45}{Hz} 的频率、\qty{2}{cm} 的振幅振动，该
\qty{40}{cm} 的摆在 $75^\circ$ 势阱中倒立，以
\qty{1.4}{Hz} 摆动。同样经时间平均的有效势能，由
振荡电四极场代替振动支点产生，
可将单个离子束缚于保罗阱中——它是原子
钟与囚禁离子量子计算的主力。
\end{solution}
